\IfFileExists{stacks-project.cls}{%
\documentclass{stacks-project}
}{%
\documentclass{amsart}
}
\usepackage{amssymb}

% For dealing with references we use the comment environment
\usepackage{verbatim}
\usepackage{fancyvrb}
\newenvironment{reference}{\comment}{\endcomment}
%\newenvironment{reference}{}{}
\newenvironment{slogan}{\comment}{\endcomment}
\newenvironment{history}{\comment}{\endcomment}

% For source figures retained as assets
\usepackage{graphicx}

% For commutative diagrams we use Xy-pic
\usepackage[all]{xy}

% We use 2cell for 2-commutative diagrams.
\xyoption{2cell}
\UseAllTwocells

% We use multicol for the list of chapters between chapters
\usepackage{multicol}

% This is generally recommended for better output
\usepackage{lmodern}
\usepackage[T1]{fontenc}
\usepackage{textalpha}

% For cross-file-references
% We prefer the snapshot tmp/xref/#2.aux,
% which latexmkrc replaces atomically, so
% that documents compiled simultaneously
% never read a half-written .aux file.
% The last argument keeps links to #2.pdf.
\usepackage{xr-hyper}
\newcommand{\maybeexternaldocument}[2]{%
  \IfFileExists{#2.aux}{%
    \externaldocument[#1]{#2}%
  }{%
    \PackageWarning{stack}{Missing #2.aux; external references with prefix #1 unavailable}%
  }%
}
\let\externaldocumentorig\externaldocument
\renewcommand{\externaldocument}[2][]{%
  \IfFileExists{tmp/xref/#2.aux}{%
    \externaldocumentorig[#1]{tmp/xref/#2}[#2.pdf]%
  }{%
  \IfFileExists{#2.aux}{%
    \externaldocumentorig[#1]{#2}%
  }{%
    \PackageWarning{stack}{Missing #2.aux; external references with prefix #1 unavailable}%
  }%
  }%
}

% Package for hypertext links:
\usepackage{hyperref}

% For any local file, say "hello.tex" you want to link to please
% use \externaldocument[hello-]{hello}
\externaldocument[introduction-]{introduction}
\externaldocument[conventions-]{conventions}
\externaldocument[sets-]{sets}
\externaldocument[categories-]{categories}
\externaldocument[topology-]{topology}
\externaldocument[sheaves-]{sheaves}
\externaldocument[sites-]{sites}
\externaldocument[stacks-]{stacks}
\externaldocument[fields-]{fields}
\externaldocument[algebra-]{algebra}
\externaldocument[brauer-]{brauer}
\externaldocument[homology-]{homology}
\externaldocument[derived-]{derived}
\externaldocument[simplicial-]{simplicial}
\externaldocument[more-algebra-]{more-algebra}
\externaldocument[smoothing-]{smoothing}
\externaldocument[modules-]{modules}
\externaldocument[sites-modules-]{sites-modules}
\externaldocument[injectives-]{injectives}
\externaldocument[cohomology-]{cohomology}
\externaldocument[sites-cohomology-]{sites-cohomology}
\externaldocument[dga-]{dga}
\externaldocument[dpa-]{dpa}
\externaldocument[sdga-]{sdga}
\externaldocument[hypercovering-]{hypercovering}
\externaldocument[schemes-]{schemes}
\externaldocument[constructions-]{constructions}
\externaldocument[properties-]{properties}
\externaldocument[morphisms-]{morphisms}
\externaldocument[coherent-]{coherent}
\externaldocument[divisors-]{divisors}
\externaldocument[limits-]{limits}
\externaldocument[varieties-]{varieties}
\externaldocument[topologies-]{topologies}
\externaldocument[descent-]{descent}
\externaldocument[perfect-]{perfect}
\externaldocument[more-morphisms-]{more-morphisms}
\externaldocument[flat-]{flat}
\externaldocument[groupoids-]{groupoids}
\externaldocument[more-groupoids-]{more-groupoids}
\externaldocument[etale-]{etale}
\externaldocument[chow-]{chow}
\externaldocument[intersection-]{intersection}
\externaldocument[pic-]{pic}
\externaldocument[weil-]{weil}
\externaldocument[adequate-]{adequate}
\externaldocument[dualizing-]{dualizing}
\externaldocument[duality-]{duality}
\externaldocument[discriminant-]{discriminant}
\externaldocument[derham-]{derham}
\externaldocument[local-cohomology-]{local-cohomology}
\externaldocument[algebraization-]{algebraization}
\externaldocument[curves-]{curves}
\externaldocument[resolve-]{resolve}
\externaldocument[models-]{models}
\externaldocument[functors-]{functors}
\externaldocument[equiv-]{equiv}
\externaldocument[pione-]{pione}
\externaldocument[etale-cohomology-]{etale-cohomology}
\externaldocument[proetale-]{proetale}
\externaldocument[relative-cycles-]{relative-cycles}
\externaldocument[more-etale-]{more-etale}
\externaldocument[trace-]{trace}
\externaldocument[crystalline-]{crystalline}
\externaldocument[spaces-]{spaces}
\externaldocument[spaces-properties-]{spaces-properties}
\externaldocument[spaces-morphisms-]{spaces-morphisms}
\externaldocument[decent-spaces-]{decent-spaces}
\externaldocument[spaces-cohomology-]{spaces-cohomology}
\externaldocument[spaces-limits-]{spaces-limits}
\externaldocument[spaces-divisors-]{spaces-divisors}
\externaldocument[spaces-over-fields-]{spaces-over-fields}
\externaldocument[spaces-topologies-]{spaces-topologies}
\externaldocument[spaces-descent-]{spaces-descent}
\externaldocument[spaces-perfect-]{spaces-perfect}
\externaldocument[spaces-more-morphisms-]{spaces-more-morphisms}
\externaldocument[spaces-flat-]{spaces-flat}
\externaldocument[spaces-groupoids-]{spaces-groupoids}
\externaldocument[spaces-more-groupoids-]{spaces-more-groupoids}
\externaldocument[bootstrap-]{bootstrap}
\externaldocument[spaces-pushouts-]{spaces-pushouts}
\externaldocument[spaces-chow-]{spaces-chow}
\externaldocument[groupoids-quotients-]{groupoids-quotients}
\externaldocument[spaces-more-cohomology-]{spaces-more-cohomology}
\externaldocument[spaces-simplicial-]{spaces-simplicial}
\externaldocument[spaces-duality-]{spaces-duality}
\externaldocument[formal-spaces-]{formal-spaces}
\externaldocument[restricted-]{restricted}
\externaldocument[spaces-resolve-]{spaces-resolve}
\externaldocument[formal-defos-]{formal-defos}
\externaldocument[defos-]{defos}
\externaldocument[cotangent-]{cotangent}
\externaldocument[examples-defos-]{examples-defos}
\externaldocument[algebraic-]{algebraic}
\externaldocument[examples-stacks-]{examples-stacks}
\externaldocument[stacks-sheaves-]{stacks-sheaves}
\externaldocument[criteria-]{criteria}
\externaldocument[artin-]{artin}
\externaldocument[quot-]{quot}
\externaldocument[stacks-properties-]{stacks-properties}
\externaldocument[stacks-morphisms-]{stacks-morphisms}
\externaldocument[stacks-limits-]{stacks-limits}
\externaldocument[stacks-cohomology-]{stacks-cohomology}
\externaldocument[stacks-perfect-]{stacks-perfect}
\externaldocument[stacks-introduction-]{stacks-introduction}
\externaldocument[stacks-more-morphisms-]{stacks-more-morphisms}
\externaldocument[stacks-geometry-]{stacks-geometry}
\externaldocument[moduli-]{moduli}
\externaldocument[moduli-curves-]{moduli-curves}
\externaldocument[examples-]{examples}
\externaldocument[exercises-]{exercises}
\externaldocument[guide-]{guide}
\externaldocument[desirables-]{desirables}
\externaldocument[coding-]{coding}
\externaldocument[obsolete-]{obsolete}
\externaldocument[fdl-]{fdl}
\externaldocument[index-]{index}

% Heap Project documents
\externaldocument[k3-]{k3}
\externaldocument[torelli-]{torelli}
\externaldocument[linear-algebra-]{linear-algebra}
\externaldocument[groups-]{groups}
\externaldocument[complex-analysis-]{complex-analysis}
\externaldocument[functional-analysis-]{functional-analysis}
\externaldocument[categories-]{categories}
\externaldocument[commutative-algebra-]{commutative-algebra}
\externaldocument[ringed-spaces-]{ringed-spaces}
\externaldocument[homological-algebra-]{homological-algebra}
\externaldocument[lie-algebras-]{lie-algebras}
\externaldocument[representation-theory-]{representation-theory}
\externaldocument[smooth-manifolds-]{smooth-manifolds}
\externaldocument[differential-topology-]{differential-topology}
\externaldocument[algebraic-topology-]{algebraic-topology}
\externaldocument[differential-geometry-]{differential-geometry}
\externaldocument[algebraic-geometry-]{algebraic-geometry}
\externaldocument[symplectic-geometry-]{symplectic-geometry}
\externaldocument[complex-geometry-]{complex-geometry}
\externaldocument[hodge-theory-]{hodge-theory}
\externaldocument[%
variations-of-hodge-structures-]{%
variations-of-hodge-structures}
\externaldocument[%
homological-mirror-symmetry-]{%
homological-mirror-symmetry}
\maybeexternaldocument{deformations-}{deformations}
\externaldocument[vertex-algebras-]{vertex-algebras}
\externaldocument[springer-]{springer}
\externaldocument[infty-categories-]{infty-categories}
\externaldocument[seminars-differential-geometry-]{%
seminars-differential-geometry}
\externaldocument[seminars-complex-geometry-]{%
seminars-complex-geometry}
\externaldocument[seminars-algebraic-geometry-]{%
seminars-algebraic-geometry}
\externaldocument[seminars-others-]{%
seminars-others}
\externaldocument[%
geometric-prequantization-]{%
geometric-prequantization}
\externaldocument[surfaces-]{surfaces}
\externaldocument[birational-maps-commutative-algebra-]{birational-maps-commutative-algebra}
\externaldocument[cremona-transformations-]{cremona-transformations}
\externaldocument[derived-categories-of-sheaves-]{derived-categories-of-sheaves}
\externaldocument[geometric-invariant-theory-]{geometric-invariant-theory}
\externaldocument[geometric-stability-conditions-and-group-actions-]{geometric-stability-conditions-and-group-actions}
\externaldocument[moduli-spaces-of-sheaves-]{moduli-spaces-of-sheaves}
\externaldocument[bridgeland-stability-general-theory-]{bridgeland-stability-general-theory}
\externaldocument[hilbert-schemes-]{hilbert-schemes}
\externaldocument[noncommutative-abelian-surfaces-and-kummer-type-hyperkahler-varieties-]{noncommutative-abelian-surfaces-and-kummer-type-hyperkahler-varieties}
\externaldocument[mumford-tate-groups-in-hodge-theory-]{mumford-tate-groups-in-hodge-theory}
\externaldocument[hodge-birational-atoms-2026-]{hodge-birational-atoms-2026}
\externaldocument[xxii-egd-2026-]{%
xxii-egd-2026}
\externaldocument[pde-]{pde}
\externaldocument[probability-]{probability}

\externaldocument[linguistics-]{linguistics}
\externaldocument[genealogia-familiar-]{%
genealogia-familiar}

\externaldocument[%
16th-alga-meeting-2026-]{%
16th-alga-meeting-2026}

\externaldocument[%
iv-jornadas-lefschetz-2026-]{%
iv-jornadas-lefschetz-2026}

\externaldocument[reading-the-masters-]{%
reading-the-masters}

\externaldocument[real-analysis-]{%
real-analysis}

\externaldocument[optimal-transport-]{%
optimal-transport}

\externaldocument[history-]{history}

\externaldocument[seiberg-witten-]{seiberg-witten}

\externaldocument[history-philosophy-biology-]{history-philosophy-biology}

\externaldocument[chemistry-]{chemistry}

\externaldocument[molecular-cell-biology-]{molecular-cell-biology}

\externaldocument[microbiology-]{microbiology}

\externaldocument[protistology-]{protistology}

\externaldocument[mycology-]{mycology}

\externaldocument[botany-]{botany}

\externaldocument[zoology-]{zoology}

\externaldocument[animal-histology-]{animal-histology}

\externaldocument[animal-physiology-]{animal-physiology}

\externaldocument[systematics-biogeography-]{systematics-biogeography}

\externaldocument[ecology-natural-resources-]{ecology-natural-resources}

\externaldocument[earth-sciences-]{earth-sciences}

\externaldocument[biostatistics-]{biostatistics}

\externaldocument[mathematical-systems-biology-]{mathematical-systems-biology}

\externaldocument[cimpa-el-salvador-2024-]{cimpa-el-salvador-2024}

% Theorem environments.
%
\theoremstyle{plain}
\newtheorem{theorem}[subsection]{Theorem}
\newtheorem{proposition}[subsection]{Proposition}
\newtheorem{lemma}[subsection]{Lemma}
\newtheorem{corollary}[subsection]{Corollary}

\theoremstyle{definition}
\newtheorem{definition}[subsection]{Definition}
\newtheorem{example}[subsection]{Example}
\newtheorem{exercise}[subsection]{Exercise}
\newtheorem{situation}[subsection]{Situation}

\theoremstyle{remark}
\newtheorem{remark}[subsection]{Remark}
\newtheorem{remarks}[subsection]{Remarks}

\numberwithin{equation}{subsection}

% Macros
%
\def\lim{\mathop{\mathrm{lim}}\nolimits}
\def\colim{\mathop{\mathrm{colim}}\nolimits}
\def\Spec{\mathop{\mathrm{Spec}}}
\def\Hom{\mathop{\mathrm{Hom}}\nolimits}
\def\Ext{\mathop{\mathrm{Ext}}\nolimits}
\def\SheafHom{\mathop{\mathcal{H}\!\mathit{om}}\nolimits}
\def\SheafExt{\mathop{\mathcal{E}\!\mathit{xt}}\nolimits}
\def\Sch{\mathit{Sch}}
\def\Mor{\mathop{\mathrm{Mor}}\nolimits}
\def\Ob{\mathop{\mathrm{Ob}}\nolimits}
\def\Sh{\mathop{\mathit{Sh}}\nolimits}
\def\NL{\mathop{N\!L}\nolimits}
\def\CH{\mathop{\mathrm{CH}}\nolimits}
\def\proetale{{pro\text{-}\acute{e}tale}}
\def\etale{{\acute{e}tale}}
\def\QCoh{\mathit{QCoh}}
\def\Ker{\mathop{\mathrm{Ker}}}
\def\Im{\mathop{\mathrm{Im}}}
\def\Coker{\mathop{\mathrm{Coker}}}
\def\Coim{\mathop{\mathrm{Coim}}}

% Boxtimes
%
\DeclareMathSymbol{\boxtimes}{\mathbin}{AMSa}{"02}

%
% Macros for moduli stacks/spaces
%
\def\QCohstack{\mathcal{QC}\!\mathit{oh}}
\def\Cohstack{\mathcal{C}\!\mathit{oh}}
\def\Spacesstack{\mathcal{S}\!\mathit{paces}}
\def\Quotfunctor{\mathrm{Quot}}
\def\Hilbfunctor{\mathrm{Hilb}}
\def\Curvesstack{\mathcal{C}\!\mathit{urves}}
\def\Polarizedstack{\mathcal{P}\!\mathit{olarized}}
\def\Complexesstack{\mathcal{C}\!\mathit{omplexes}}
% \Pic is the operator that assigns to X its picard group, usage \Pic(X)
% \Picardstack_{X/B} denotes the Picard stack of X over B
% \Picardfunctor_{X/B} denotes the Picard functor of X over B
\def\Pic{\mathop{\mathrm{Pic}}\nolimits}
\def\Picardstack{\mathcal{P}\!\mathit{ic}}
\def\Picardfunctor{\mathrm{Pic}}
\def\Deformationcategory{\mathcal{D}\!\mathit{ef}}


\begin{document}

\title{Molecular and Cell Biology}
\maketitle

\phantomsection
\label{section-phantom}

\tableofcontents

\section{What this chapter covers}
\label{section-mcb-overview}

These notes collect the molecular and cell
biology of Daniel's first three undergraduate
years: nucleic acids and the flow of genetic
information, enzymes, hemoglobin, membranes
and energy, the cell biology of junctions,
the cell cycle, signalling and death, the
laboratory techniques he used, and the
biotechnology built on them.

\textbf{Sources of these notes.}
Three courses of the Facultad de Ciencias,
UNAM. Biología Molecular de la Célula I (2nd
semester, 2014-2): his answers to the
departmental enzymology exam, his study
questionnaire on nucleic acids, a talk on the
molecular basis of hemoglobinopathies (given
with a classmate), a laboratory report on
protein quantification, a theoretical cloning
problem with the plasmid pQE30, and his
outline for the laboratory exam. Biología
Molecular de la Célula III (4th semester,
2015-2): twelve short practice reports on
microscopy and staining, two exam outlines
(junctions, extracellular matrix, cell cycle
and meiosis; signalling, differentiation and
cell death) and two sets of lecture slides.
Biotecnología (6th semester, 2016-2): lecture
slides on history, genes and genomes,
transcription and PCR, the course manual, and
his report on DNA extraction from soy sprouts
for GMO detection. The material was checked
and completed with Alberts et al.,
\emph{Essential Cell Biology}
\cite{alberts2014essential}, Voet, Voet and
Pratt, \emph{Fundamentals of Biochemistry}
\cite{voet2013fundamentals}, and the methods
chapter of Ross and Pawlina
\cite{ross2011histology}.

\section{Nucleic acids and genomes}
\label{section-mcb-nucleic-acids}

A \emph{nucleoside} is a base joined to a
pentose (ribose in RNA, 2'-deoxyribose in
DNA) by a $\beta$-N-glycosidic bond at C1'; a
\emph{nucleotide} is a nucleoside
phosphorylated at C5'. The purines are
adenine and guanine (attached through N9);
the pyrimidines are cytosine, thymine and
uracil (attached through N1). Nucleotides are
joined by 3'-5' phosphodiester bonds, so
every strand has a polarity, written from the
5' end to the 3' end \cite[Section
3-2]{voet2013fundamentals}.

\textbf{The double helix.}
Chargaff found, in the late 1940s, that DNA
has as much A as T and as much G as C, while
the G+C content varies between species
(25--75\% in bacteria). Watson and Crick
explained this by base pairing in an
antiparallel double helix: A pairs with T by
two hydrogen bonds, G with C by three. In the
B form the bases sit inside, nearly
perpendicular to the axis, the
sugar-phosphate backbones wind outside and
form a major and a minor groove, and the
helix is 20 \AA{} wide \cite[Sections 3-2 and
24-1]{voet2013fundamentals}.

\begin{center}
\begin{tabular}{llll}
Form & Hand & bp/turn & Notes \\
\hline
B & right & about 10 & usual form \\
A & right & 11.6 & dehydrated; \\
 & & & bases tilted $20^\circ$ \\
Z & left & 12 & alternating \\
 & & & purine-pyrimidine \\
\end{tabular}
\end{center}

The A form has a deep major groove, a shallow
minor groove and an axial hole; the Z form
(found in d(CGCGCG) crystals, favoured by
high salt) has a zig-zag backbone and is
probably present in vivo, since Z-DNA-binding
proteins exist \cite[Section
24-1]{voet2013fundamentals}. Meselson and
Stahl showed with $^{15}$N density labelling
that replication is semiconservative
\cite[Chapter 6]{alberts2014essential}.

\textbf{Genome sizes.}
From the table of sequenced genomes in
\cite[Section 3-4]{voet2013fundamentals}, in
kilobases: \emph{Escherichia coli} 4,639;
\emph{Saccharomyces cerevisiae} 12,070;
\emph{Drosophila melanogaster} 180,000; rice
389,000; \emph{Homo sapiens} 3,038,000 (about
$3 \times 10^9$ bp in 23 chromosomes). The
human genome has roughly 20,000--25,000
protein-coding genes; rice has more. Neither
genome size nor gene number tracks complexity
in a simple way (the C-value and G-value
paradoxes of the biotechnology course;
\cite[Chapter 5]{alberts2014essential}).

\textbf{RNAs.}
Only mRNA is translated. Non-coding RNAs
include rRNA (ribosome structure and peptidyl
transfer), tRNA (adaptor), snRNA (splicing),
and the small regulatory miRNA and siRNA
(18--25 nucleotides, silencing complementary
mRNAs) \cite[Chapter
28]{voet2013fundamentals}.

\section{DNA replication}
\label{section-mcb-replication}

Replication starts at origins (one,
\emph{oriC}, in \emph{E. coli}; many in
eukaryotes), proceeds in both directions, and
is always synthesized 5' to 3'. One strand
(leading) is made continuously; the other
(lagging) in Okazaki fragments, each started
on a short RNA primer made by primase and
later joined by ligase \cite[Chapter
6]{alberts2014essential}.

\begin{center}
\begin{tabular}{lll}
Function & Bacteria & Eukaryotes \\
\hline
main polymerase & Pol III & Pol $\varepsilon$
(leading), \\
 & & Pol $\delta$ (lagging) \\
primer & primase & Pol $\alpha$-primase \\
primer removal & Pol I & RNase H, FEN1 \\
helicase & DnaB & MCM \\
single strands & SSB & RPA \\
supercoils & gyrase & topoisomerases \\
 & (topo II) & I and II \\
origin & DnaA at & ORC \\
recognition & oriC & \\
chromosome ends & (circular) & telomerase \\
\end{tabular}
\end{center}

Pol $\gamma$ replicates mitochondrial DNA.
Telomerase carries its own RNA template and
extends the 3' ends that the lagging strand
machinery cannot finish. In eukaryotes the
MCM helicase is loaded on ORC-bound origins
in G1 (licensing) and activated in S phase
only, so each origin fires once per cycle
\cite[Chapter 25]{voet2013fundamentals}.

\section{Transcription, translation and the
genetic code}
\label{section-mcb-expression}

\textbf{Bacteria.}
One RNA polymerase (core
$\alpha_2\beta\beta'\omega$ plus a $\sigma$
factor that recognizes the promoter) makes
all RNAs. A typical promoter has the
consensus TATAAT at $-10$ (Pribnow box) and
TTGACA at $-35$. The polymerase opens a
bubble of about 17 bp; $\sigma$ is released
after the first few nucleotides. Termination
uses either a GC-rich hairpin followed by a
run of U, or the protein rho \cite[Chapter
26]{voet2013fundamentals}.

\textbf{Eukaryotes.}
RNA polymerase I makes the large rRNA
precursor, II the pre-mRNAs (and most
snRNAs), III the tRNAs and 5S rRNA.
Polymerase II needs general transcription
factors: TFIID (with the TATA-binding
protein, TBP, on the TATA box 25--31 bp
upstream), TFIIB, TFIIE, TFIIF and TFIIH (a
helicase that opens the DNA and
phosphorylates the C-terminal domain of the
polymerase, which releases it into
elongation). The Mediator complex relays
signals from gene-specific activators. A
CCAAT box often lies further upstream (around
$-70$ to $-90$). The primary transcript is
capped at the 5' end (7-methylguanosine),
cleaved 15--25 nucleotides past an AAUAAA
signal and given a poly(A) tail, and spliced
by the spliceosome (snRNPs U1, U2, U4, U5,
U6), which removes introns as lariats
\cite[Chapter 26]{voet2013fundamentals}.

\textbf{Translation.}
Aminoacyl-tRNA synthetases load each tRNA
with its amino acid,
$$
\mathrm{aa} + \mathrm{tRNA} + \mathrm{ATP}
\rightarrow \mathrm{aa{-}tRNA} + \mathrm{AMP}
+ \mathrm{PP_i},
$$
and fall into two unrelated classes, I and
II, of ten enzymes each \cite[Chapter
27]{voet2013fundamentals}.

\begin{center}
\begin{tabular}{lll}
Step & Bacteria & Eukaryotes \\
\hline
start signal & Shine-Dalgarno & 5' cap; scan
\\
 & sequence & to first AUG \\
initiator & fMet-tRNA & Met-tRNA \\
initiation & IF1, IF2, IF3 & a dozen or \\
factors & & more eIFs \\
elongation & EF-Tu, EF-Ts & eEF1A, eEF1B \\
translocation & EF-G & eEF2 \\
release & RF1, RF2, RF3 & eRF1 (all \\
 & & three stops) \\
recycling & RRF + EF-G & \\
\end{tabular}
\end{center}

The Shine-Dalgarno sequence is a purine-rich
tract about 10 nucleotides upstream of the
start codon that pairs with the 3' end of 16S
rRNA \cite[Chapter 27]{voet2013fundamentals}.

\textbf{The code.}
With four bases, triplets are the shortest
words that can name 20 amino acids: $4^3 =
64$ codons, 61 for amino acids and three
stops (UAA, UAG, UGA). The code is
\emph{degenerate} (most amino acids have
several codons) and AUG is the usual start.
Nirenberg and Matthaei found in 1961 that
poly(U) directs the synthesis of
polyphenylalanine, so UUU codes for Phe
\cite[Section 27-1]{voet2013fundamentals}. A
DNA sequence has six reading frames (three
per strand); an open reading frame runs from
a start codon to a stop codon.

\section{Gene regulation}
\label{section-mcb-regulation}

In bacteria, genes of one pathway are often
grouped in an operon read as one mRNA. The
\emph{lac} operon is shut by the Lac
repressor when lactose is absent and needs
the activator CAP (with cAMP, abundant when
glucose is scarce) to be transcribed
strongly; the \emph{trp} operon is repressed
when tryptophan is abundant \cite[Chapter
8]{alberts2014essential}.

In eukaryotes, control is mostly at
transcription initiation: gene-specific
activators and repressors bind enhancers and
silencers, often thousands of base pairs
away, and act through Mediator and through
chromatin. DNA is wrapped around histone
octamers (H2A, H2B, H3, H4, rich in Lys and
Arg) into nucleosomes, sealed by H1; open
euchromatin is transcribed, condensed
heterochromatin mostly not. Later controls
include alternative splicing, mRNA stability
and silencing by miRNAs \cite[Chapter
8]{alberts2014essential}.

\section{Enzymes}
\label{section-mcb-enzymes}

An enzyme lowers the activation free energy
$\Delta G^{\ddagger}$ of a reaction by
stabilizing its transition state. It changes
neither $\Delta G$ nor the equilibrium
constant $K$, which are fixed by the
reactants and products, so it speeds up both
directions equally \cite[Section
11-2]{voet2013fundamentals}. Fischer's
lock-and-key model (1894) explained
specificity by geometric and electronic
complementarity between enzyme and substrate;
Koshland's induced fit adds that substrate
binding changes the enzyme's conformation,
which is also the basis of his sequential
model of allostery \cite[Sections 7-1 and
11-1]{voet2013fundamentals}.

\textbf{Catalytic mechanisms.}
Voet classifies them into five types
\cite[Section 11-3]{voet2013fundamentals}:
\begin{enumerate}
\item acid-base catalysis: proton transfer to
or from side chains such as His, Asp, Glu
(RNase A);
\item covalent catalysis: a nucleophile of
the enzyme (Ser, Cys, His) forms a transient
covalent intermediate;
\item metal ion catalysis: bound metal ions
orient the substrate, stabilize charges or
change oxidation state;
\item proximity and orientation effects (with
electrostatic catalysis as a variant);
\item preferential binding of the transition
state, which is bound more tightly than the
substrate.
\end{enumerate}

\textbf{Michaelis-Menten kinetics.}
For $\mathrm{E} + \mathrm{S}
\rightleftharpoons \mathrm{ES} \rightarrow
\mathrm{E} + \mathrm{P}$ with $\mathrm{ES}$
at steady state, the initial velocity is
\begin{equation}
\label{equation-mcb-michaelis-menten}
v = \frac{V_{\max}[\mathrm{S}]}{K_M +
[\mathrm{S}]}.
\end{equation}
$K_M$ is the substrate concentration giving
$v = V_{\max}/2$. Taking reciprocals gives
the Lineweaver-Burk (double-reciprocal) line
\begin{equation}
\label{equation-mcb-lineweaver-burk}
\frac{1}{v} = \frac{K_M}{V_{\max}} \cdot
\frac{1}{[\mathrm{S}]} + \frac{1}{V_{\max}},
\end{equation}
with intercept $1/V_{\max}$ on the $1/v$ axis
and $-1/K_M$ on the $1/[\mathrm{S}]$ axis
\cite[Section 12-1]{voet2013fundamentals}.

\begin{example}
\label{example-mcb-exam-kinetics}
The departmental exam gave these initial
velocities (in nmol L$^{-1}$ min$^{-1}$) for
a Michaelis-Menten enzyme:

\begin{center}
\begin{tabular}{rr}
$[\mathrm{S}]$ (M) & $v$ \\
\hline
$6.25 \times 10^{-6}$ & 15.0 \\
$7.50 \times 10^{-5}$ & 56.25 \\
$1.00 \times 10^{-4}$ & 60.0 \\
$1.00 \times 10^{-3}$ & 74.9 \\
$1.00 \times 10^{-2}$ & 75 \\
\end{tabular}
\end{center}

A least-squares line through the five points
$(1/[\mathrm{S}], 1/v)$ gives (recomputed for
these notes)
$$
\frac{1}{v} = 3.340 \times 10^{-7} \cdot
\frac{1}{[\mathrm{S}]} + 0.013239,
$$
so $V_{\max} = 1/0.013239 = 75.5$ nmol
L$^{-1}$ min$^{-1}$ and $K_M = 3.340 \times
10^{-7} \times 75.5 = 2.52 \times 10^{-5}$ M,
in agreement with Daniel's answer. The data
were evidently built from $V_{\max} = 75$ and
$K_M = 2.5 \times 10^{-5}$ M: these values
reproduce the first three velocities exactly.
The two velocities asked for then follow from
\eqref{equation-mcb-michaelis-menten}: at
$[\mathrm{S}] = 2.5 \times 10^{-5}$ M, which
equals $K_M$, $v = V_{\max}/2 \approx 37.6$;
at $[\mathrm{S}] = 5.0 \times 10^{-5}$ M, $v
= 2V_{\max}/3 \approx 50.2$ nmol L$^{-1}$
min$^{-1}$.
\end{example}

\textbf{Protein folding.}
The amino acid sequence contains the
information for the native structure
(Anfinsen's renaturation of ribonuclease),
but folding in the cell follows pathways
through partly folded intermediates and is
assisted by molecular chaperones; misfolding
and aggregation cause diseases such as the
amyloidoses \cite[Section
6-5]{voet2013fundamentals}. The course's
reading on this was Gómez Puyou's review on
folding in vitro and in vivo (Mensaje
Bioquímico, 2012).

\section{Hemoglobin and hemoglobinopathies}
\label{section-mcb-hemoglobin}

Hemoglobin is a tetramer of two $\alpha$-like
and two $\beta$-like globins, each carrying a
heme whose Fe(II) binds $\mathrm{O_2}$.
Oxygen binding is cooperative (Hill
coefficient 2.8--3.0, sigmoidal curve,
half-saturation at 26 torr in blood) because
the tetramer switches from the low-affinity T
state to the high-affinity R state.
2,3-Bisphosphoglycerate (BPG) binds in the
central cavity of deoxyhemoglobin, stabilizes
T and lowers the $\mathrm{O_2}$ affinity,
which lets hemoglobin unload oxygen in the
tissues \cite[Section
7-1]{voet2013fundamentals}.

\textbf{Globin genes and switching.}
The $\alpha$-like genes ($\zeta$, $\alpha_1$,
$\alpha_2$) lie in a cluster on chromosome
16; the $\beta$-like genes ($\varepsilon$,
$^G\gamma$, $^A\gamma$, $\delta$, $\beta$) on
chromosome 11, in the order of their
expression in development \cite[Chapter
28]{voet2013fundamentals}.

\begin{center}
\begin{tabular}{lll}
Stage & Hemoglobin & Chains \\
\hline
embryo & Gower 1 & $\zeta_2\varepsilon_2$ \\
 & Gower 2 & $\alpha_2\varepsilon_2$ \\
 & Portland & $\zeta_2\gamma_2$ \\
fetus & HbF & $\alpha_2\gamma_2$ \\
adult & HbA (97\%) & $\alpha_2\beta_2$ \\
 & HbA$_2$ (2\%) & $\alpha_2\delta_2$ \\
 & HbF (1\%) & $\alpha_2\gamma_2$ \\
\end{tabular}
\end{center}

HbF binds BPG more weakly than HbA, so it has
a higher $\mathrm{O_2}$ affinity and takes
oxygen from the maternal blood across the
placenta \cite[Section
7-1]{voet2013fundamentals}.

\textbf{Structural variants.}
Over 1000 variants are known, most from
single amino acid substitutions.
Substitutions on the surface are usually
harmless; those inside the molecule cause
disease \cite[Section
7-1]{voet2013fundamentals}:
\begin{itemize}
\item unstable hemoglobins (heme binding or
packing weakened) precipitate as Heinz bodies
and cause hemolytic anemia;
\item methemoglobins (HbM; substitutions near
the heme, e.g. Hb Boston, His E7 to Tyr)
favour Fe(III), which cannot bind
$\mathrm{O_2}$: heterozygotes are cyanotic,
and no homozygotes are known;
\item altered affinity: Hb Kansas ($\beta$102
Asn to Thr) loses a hydrogen bond of the R
state and has low affinity (mild cyanosis);
Hb Yakima ($\beta$99 Asp to His) loses a
hydrogen bond of the T state and has high
affinity, which causes polycythemia; Hb
Kempsey ($\beta$99 Asp to Asn), the talk's
example, is a high-affinity variant at the
same position.
\end{itemize}

\textbf{Sickle-cell anemia.}
In HbS, Glu $\beta$6 is replaced by Val
(Ingram, 1956). In the deoxy state the Val
fits a hydrophobic pocket (the talk named its
Phe and Leu residues) on a $\beta$ subunit of
a neighbouring tetramer; HbS polymerizes into
fibres of 14 strands that deform the red
cells into sickles. Oxygenation dissolves the
fibres, so sickling happens in capillaries.
Homozygotes ($\beta^S\beta^S$) have chronic
hemolytic anemia and painful vaso-occlusive
crises; heterozygotes (sickle trait, about
40\% HbS) are usually healthy and are more
resistant to \emph{Plasmodium falciparum}
malaria, which keeps the allele frequent
where malaria is endemic. High HbF dilutes
HbS and softens the disease; hydroxyurea
raises HbF \cite[Section
7-1]{voet2013fundamentals}.

\textbf{Thalassemias.}
These are defects in the amount, not the
structure, of a globin chain. Many
$\beta$-thalassemias are point mutations that
affect transcription, splicing or
translation; in one patient a G to A change
in the first intron created a new 3' splice
site and a premature stop \cite[Section
28.3.3]{bergbiochemistry}. The rest of
this paragraph follows Daniel's talk and was
not checked against the textbooks.
$\alpha$-Thalassemias are mostly deletions of
one to all four $\alpha$ genes (or of the
upstream regulatory element HS-40); the
excess chains form Hb Bart's ($\gamma_4$) in
the fetus and HbH ($\beta_4$) in adults.
$\beta$-Thalassemias matter after birth, when
$\beta$ replaces $\gamma$, and raise HbA$_2$
and HbF. Unequal crossing-over between
$\delta$ and $\beta$ makes the fusion chain
of Hb Lepore ($\alpha_2(\delta\beta)_2$). HbE
($\beta$26 Glu to Lys) is a structural
variant with a thalassemic phenotype.

\textbf{Diagnosis and treatment (from the
talk).}
Blood smears, hemoglobin electrophoresis and
spectroscopy for diagnosis; transfusion (with
iron overload as a risk), splenectomy, bone
marrow transplantation, and gene therapy,
which by 2013 had treated $\beta$-thalassemia
\cite[Section 3-5]{voet2013fundamentals}.

\section{Membranes and cell metabolism}
\label{section-mcb-membranes-metabolism}

This section is written from the textbooks;
the surviving coursework has no material on
it.

\textbf{The lipid bilayer.}
Membrane lipids (phospholipids, cholesterol,
glycolipids) are amphipathic: a hydrophilic
head and hydrophobic tails. In water they
assemble spontaneously into a bilayer, which
is a two-dimensional fluid; its fluidity
rises with shorter and unsaturated (kinked)
tails and is buffered by cholesterol. The two
leaflets differ in composition
(phosphatidylserine faces the cytosol;
glycolipids face outside). Proteins are
embedded in or attached to the bilayer and
many drift laterally \cite[Chapter
11]{alberts2014essential}.

\textbf{Transport.}
Small nonpolar molecules ($\mathrm{O_2}$,
$\mathrm{CO_2}$) diffuse through the bilayer;
ions and polar molecules need transport
proteins \cite[Chapter
12]{alberts2014essential}:
\begin{itemize}
\item channels (ion channels, aquaporins) let
solutes pass down their electrochemical
gradient, fast and selectively;
\item transporters bind the solute and change
conformation: passively (glucose uniporter)
or actively, using ATP (the
$\mathrm{Na^+}$-$\mathrm{K^+}$ pump, 3
$\mathrm{Na^+}$ out and 2 $\mathrm{K^+}$ in
per ATP) or the downhill flow of another ion
(symporters and antiporters coupled to
$\mathrm{Na^+}$ or $\mathrm{H^+}$);
\item water follows solutes by osmosis.
\end{itemize}
The pump and the $\mathrm{K^+}$ leak channels
set the resting membrane potential, inside
negative.

\textbf{Energy from food.}
ATP is the energy currency: its hydrolysis is
strongly favourable and is coupled to
unfavourable reactions; NADH and FADH$_2$
carry high-energy electrons. Glycolysis, in
the cytosol, splits glucose into two
pyruvates, with a net gain of 2 ATP and 2
NADH per glucose. Without oxygen, pyruvate is
reduced to lactate or ethanol (fermentation)
to regenerate NAD$^+$. With oxygen, pyruvate
enters the mitochondrion, becomes acetyl CoA,
and is oxidized to $\mathrm{CO_2}$ in the
citric acid cycle, which yields NADH,
FADH$_2$ and GTP \cite[Chapter
13]{alberts2014essential}.

\textbf{Chemiosmotic coupling.}
In the inner mitochondrial membrane, the
electron transport chain passes electrons
from NADH to $\mathrm{O_2}$ and uses the
energy to pump protons out of the matrix. The
resulting electrochemical proton gradient
drives protons back through ATP synthase,
which makes ATP (Mitchell's chemiosmotic
mechanism); the same gradient drives other
transport. Complete oxidation of one glucose
gives about 30 ATP, against 2 from glycolysis
alone. Chloroplasts use the same principle:
light-driven electron transport in the
thylakoid membrane pumps protons and makes
ATP and NADPH for carbon fixation
\cite[Chapter 14]{alberts2014essential}.

\section{Cell junctions and the extracellular
matrix}
\label{section-mcb-junctions}

The lecture classification of junctions
\cite[Chapter 20]{alberts2014essential}:
\begin{itemize}
\item anchoring junctions, linked to the
cytoskeleton: adherens junctions (cadherin to
actin, cell-cell), desmosomes (cadherin to
intermediate filaments, cell-cell),
actin-linked cell-matrix junctions (integrin
to actin) and hemidesmosomes (integrin to
intermediate filaments, cell to basal
lamina);
\item occluding junctions: tight junctions in
vertebrates (claudins and occludins), which
seal epithelia and keep membrane proteins in
their apical or basolateral domain; septate
junctions in invertebrates;
\item channel-forming junctions: gap
junctions (connexins in vertebrates, innexins
in invertebrates; ions and small metabolites)
and plasmodesmata in plants;
\item signal-relaying junctions, such as
chemical synapses.
\end{itemize}
Along an epithelium the order from the apex
is tight junction, adherens junction,
desmosome, with gap junctions further down
and hemidesmosomes at the base.

Cadherins are $\mathrm{Ca^{2+}}$-dependent,
mostly homophilic (a cadherin binds the same
cadherin on the next cell), which lets cells
sort into tissues. Integrins bind matrix
proteins. Selectins bind carbohydrates and
mediate the transient adhesion of leukocytes
to the endothelium.

The extracellular matrix of connective tissue
is made of glycosaminoglycans (mostly bound
to proteins as proteoglycans), which form a
hydrated gel resisting compression, and
fibrous proteins: collagen (a triple helix
with Gly at every third residue, rich in Pro;
about a quarter of mammalian protein) for
tensile strength, elastin for elasticity, and
the adhesive proteins fibronectin and laminin
(the latter in the basal lamina)
\cite[Chapter 20]{alberts2014essential}.

\section{The cell cycle, mitosis and meiosis}
\label{section-mcb-cell-cycle}

The cycle is G1, S (DNA replication), G2 and
M (mitosis and cytokinesis); G1, S and G2
form interphase. In cultured human cells the
lecture gave about 11 h, 8 h, 4 h and 1 h,
about 24 h in total. Cells that stop dividing
leave G1 for G0. Passing the restriction
point (Start in yeast) at the end of G1
commits the cell to divide \cite[Chapter
18]{alberts2014essential}.

\textbf{Control.}
The engine is a set of cyclin-dependent
kinases (Cdks), active only when bound to a
cyclin; cyclin levels rise and fall through
synthesis and degradation. MPF, found as an
activity in frog eggs that drives oocytes
into M phase, is M-Cdk: Cdk1 (Cdc2) with
cyclin B. Cdk4/6-cyclin D acts in G1,
Cdk2-cyclin E at G1/S, Cdk2-cyclin A in S.
Checkpoints stop the cycle if DNA is damaged
or unreplicated or if chromosomes are not
attached to the spindle: DNA damage
stabilizes p53, which induces the Cdk
inhibitor p21; if the damage cannot be
repaired, p53 triggers apoptosis
\cite[Chapter 18]{alberts2014essential}. The
yeasts \emph{S. pombe} and \emph{S.
cerevisiae}, with temperature-sensitive
\emph{cdc} mutants, were the models; normal
human cells divide only a limited number of
times in culture.

\textbf{Mitosis.}
Prophase (chromosomes condense; the spindle
forms between the duplicated centrosomes),
prometaphase (the nuclear envelope breaks
down; kinetochores attach to spindle
microtubules), metaphase (chromosomes align
at the equator), anaphase (sister chromatids
separate and move to the poles), telophase
(nuclear envelopes reform) and cytokinesis by
an actin-myosin contractile ring
\cite[Chapter 18]{alberts2014essential}. In
the course's onion root tip squash stained
with acetic orcein ($2n = 16$ for
\emph{Allium cepa}), all these stages were
seen.

\textbf{Meiosis.}
One round of replication followed by two
divisions gives four haploid cells. In
prophase I homologues pair into bivalents and
cross over \cite[Chapter
19]{alberts2014essential}. Its stages are
leptotene (chromatin condenses), zygotene
(homologues pair through the synaptonemal
complex), pachytene (crossing-over),
diplotene (the synaptonemal complex
dissolves; homologues stay joined at
chiasmata) and diakinesis (further
condensation) \cite[Chapter
3]{ross2011histology}.

\section{Cell signalling, differentiation and
cell death}
\label{section-mcb-signalling}

\textbf{Signalling.}
Signals act at different ranges: endocrine
(hormones in the blood), paracrine (local
mediators), neuronal (synapses) and
contact-dependent. A signal is synthesized,
secreted, transported, received by a
receptor, relayed by an intracellular
cascade, and finally removed; the response
may change enzyme activity, the cytoskeleton,
ion permeability, gene expression or survival
\cite[Chapter 16]{alberts2014essential}.

G-protein-coupled receptors cross the
membrane seven times. The trimeric G protein
carries GDP on its $\alpha$ subunit when
inactive; the activated receptor makes
$\alpha$ exchange GDP for GTP, and $\alpha$
and $\beta\gamma$ separate and act on
targets. $\mathrm{G_s}$ $\alpha$ activates
adenylyl cyclase, which makes the second
messenger cAMP from ATP; cAMP activates
protein kinase A. The signal ends when
$\alpha$ hydrolyses its GTP and when the
receptor is phosphorylated and bound by
arrestin \cite[Chapter
16]{alberts2014essential}.

\textbf{Differentiation.}
Cells become different through differences in
gene expression, set by cytoplasmic
determinants, signals from neighbours and the
matrix. By potential: totipotent (the zygote;
all cells including extraembryonic ones),
pluripotent (all cells of the body, e.g.
embryonic stem cells), multipotent (several
types, e.g. hematopoietic stem cells),
unipotent. Homeotic (Hox) genes code for
transcription factors with a DNA-binding
homeodomain, arranged on the chromosome in
the order of the body regions they specify
\cite[Chapter 20]{alberts2014essential}.

\textbf{Cell death.}
Three forms were contrasted in the course:
\begin{itemize}
\item necrosis: accidental; the cell swells
and bursts, spilling its contents and causing
inflammation;
\item apoptosis: programmed and
ATP-dependent; the cell shrinks, chromatin
condenses, the nucleus and DNA fragment
(between nucleosomes), and the cell breaks
into membrane-bounded apoptotic bodies that
are engulfed. It is executed by caspases
(cysteine proteases). In the intrinsic
pathway, Bax and Bak make the outer
mitochondrial membrane leaky and cytochrome c
(with Smac/DIABLO, which blocks caspase
inhibitors) is released; in the extrinsic
pathway death receptors activate caspases.
Exposed phosphatidylserine, detected with
annexin V, marks apoptotic cells
\cite[Chapter 18]{alberts2014essential};
\item autophagy: the cell digests its own
contents in autophagosomes fused with
lysosomes, typically under starvation; it
does not necessarily kill the cell (Atg
genes; the lysosomal marker LAMP).
\end{itemize}

\section{Microscopy and staining}
\label{section-mcb-microscopy}

\textbf{Köhler illumination.}
It makes the illuminated field exactly match
the field of view, evenly lit. The steps in
Ross \cite[Chapter 1]{ross2011histology}:
focus the specimen; close the field
diaphragm; move the condenser until the edge
of the field diaphragm (the polygon seen in
the course) is sharp; centre it; open it
until it just clears the field; then set the
condenser (aperture) diaphragm to about two
thirds of the objective aperture, the best
compromise between resolution and contrast.

\textbf{Phase contrast.}
Unstained cells differ from their medium in
refractive index, which shifts the phase of
light but barely its intensity. A phase ring
in the objective and an annulus in the
condenser turn these phase shifts into
differences in brightness \cite[Chapter
1]{ross2011histology}. In the course,
unstained cheek epithelial cells were nearly
invisible in bright field and clear in phase
contrast; \emph{Elodea} leaves showed
chloroplasts carried around the cell by
cytoplasmic streaming (cyclosis).

\textbf{Fluorescence.}
The fluorochrome is excited by light of one
wavelength and emits at a longer one; an
excitation filter selects the exciting light
and an emission filter passes only the
emitted light. DAPI binds the minor groove of
AT-rich DNA and, when bound, absorbs at 358
nm (ultraviolet) and emits at 461 nm (blue),
so it lights up nuclei.

\textbf{Preparing and staining tissue.}
Fixation, dehydration through graded
ethanols, clearing in xylene and embedding in
paraffin give blocks cut at a few micrometres
(5--20 $\mu$m in the course notes); semithin
(about 0.5--1 $\mu$m) and ultrathin (60--90
nm, for electron microscopy) sections need
resin. To stain, sections are deparaffinized
in xylene and rehydrated through descending
ethanols to water. Basic dyes stain acidic
(basophilic) components such as nucleic
acids; acidic dyes stain basic (acidophilic)
components such as most cytoplasmic proteins
\cite[Chapter 1]{ross2011histology}. The
stains used in the course:
\begin{itemize}
\item hematoxylin (used with a mordant, which
makes it behave as a basic dye) and eosin
(acidic): nuclei blue-purple, cytoplasm pink;
used on testis;
\item toluidine blue (basic, metachromatic):
rat ovary follicles;
\item Giemsa on methanol-fixed blood smears;
\item methyl green (basic): nuclei;
\item silver impregnation (Golgi's method;
the AgNOR technique stains nucleolar
organizer regions): cerebellum;
\item acetic orcein: chromosomes in onion
root tips.
\end{itemize}

\section{Molecular biology laboratory
techniques}
\label{section-mcb-lab-techniques}

\textbf{Bradford assay.}
Coomassie Brilliant Blue G-250 binds proteins
(mainly basic and aromatic residues) and its
absorbance maximum shifts to 595 nm, in
proportion to the protein present. Unknowns
are read against a standard curve. In the
course's practice, casein was first isolated
from skimmed milk powder by isoelectric
precipitation; the standard curve used 1 to
15 $\mu$g of bovine serum albumin with 400
$\mu$L of reagent in 2 mL. The averaged
curve, $A_{595} = 0.0213 m - 0.0013$ with $m$
in $\mu$g, gave about 0.15 g of casein per
gram of powder, against about 0.36 g of total
protein per gram on the label.

\textbf{Transformation and plasmid
preparation.}
Plasmids are small circular DNAs with their
own origin of replication (copy number from
one to hundreds per cell), a selectable
marker and, in vectors, a polylinker of
unique restriction sites. Bacteria made
competent with divalent cations
($\mathrm{CaCl_2}$) take up plasmid after a
heat shock (42 $^\circ$C, 45 s), are allowed
to express the resistance gene in LB, and are
plated on ampicillin; the \emph{bla} gene
encodes a $\beta$-lactamase that hydrolyses
the $\beta$-lactam ring. Plasmid is recovered
by alkaline lysis: resuspension (Tris, EDTA,
glucose), lysis with NaOH and SDS (DNA
denatures), neutralization with potassium or
sodium acetate (chromosomal DNA and proteins
precipitate with the SDS, the small plasmid
renatures), and ethanol precipitation of the
supernatant. pGLO, used in the course,
carries the green fluorescent protein gene
under an arabinose-inducible promoter, so
transformed colonies glow on arabinose
plates.

\textbf{Gels.}
DNA is negatively charged and runs to the
anode; in agarose (typically 1\%) it
separates by length, the distance migrated
decreasing roughly linearly with the
logarithm of the size, so sizes are read
against a ladder \cite[Chapter
10]{alberts2014essential}. Small fragments
are resolved better in polyacrylamide. In
SDS-PAGE, proteins are denatured with SDS and
a reducing agent (Laemmli buffer); SDS coats
them with a uniform negative charge, so they
separate by mass and are seen after Coomassie
staining \cite[Section
5-2]{voet2013fundamentals}.

\textbf{Restriction mapping.}
Type II restriction enzymes cut at short
palindromic sites, e.g. EcoRI cuts G/AATTC
and HindIII A/AGCTT, leaving sticky ends
\cite[Section 3-4]{voet2013fundamentals}.

\begin{example}
\label{example-mcb-pqe30}
The extra problem of the first course: prove
that the \emph{ras} gene of the dab
\emph{Limanda limanda} has been cloned into
the EcoRI site of the expression vector
pQE30, starting from 1 $\mu$L of plasmid.
Daniel's design: transform \emph{E. coli}
HB101 (with a no-plasmid control), grow on LB
with ampicillin, prepare plasmid by alkaline
lysis, digest with EcoRI, HindIII and both,
and run a 1\% agarose gel next to a 1 kb
ladder. The given sequence has 4,636 bp,
which is pQE30 (3,461 bp) plus the 1,175 bp
insert; it has one EcoRI site (position 88)
and one HindIII site (position 1,362), as was
checked again for these notes. Each single
digest therefore gives one linear band of
4,636 bp, and the double digest gives two
bands, $1,362 - 88 = 1,274$ bp and $4,636 -
1,274 = 3,362$ bp. A plasmid without the
insert would give 3,461 bp in the single
digests, so the bands prove the insertion. By
translating the open reading frames, he also
identified the vector's $\beta$-lactamase.
\end{example}

\section{Biotechnology and recombinant DNA}
\label{section-mcb-biotechnology}

Biotechnology uses organisms, or their parts,
to make products or provide services: in
health, agriculture, animals, industry and
the environment. Its empirical phase
(fermented drinks and bread, selective
breeding) is thousands of years old; the
history of winemaking and yeast research was
one of the course readings (Chambers and
Pretorius, EMBO Reports 2010).

\begin{center}
\begin{tabular}{ll}
Year & Milestone \\
\hline
1856--1860s & Pasteur: fermentation \\
 & is done by microbes \\
1897 & Buchner: fermentation \\
 & in cell-free yeast extract \\
1919 & Ereky coins \\
 & \emph{Biotechnologie} \\
1944 & Avery, MacLeod, McCarty: \\
 & DNA is the transforming \\
 & principle \\
1953 & double helix \\
1958 & Meselson-Stahl; \\
 & Kornberg's DNA polymerase \\
1961--1966 & genetic code deciphered \\
1972 & Berg: first recombinant \\
 & DNA molecule \\
1973 & Cohen and Boyer clone \\
 & a plasmid in bacteria \\
1977 & Sanger sequencing \\
1982 & recombinant human insulin \\
 & approved \\
1983--1985 & Mullis: PCR \\
1996 & first large-scale GM crops \\
2001--2004 & human genome draft \\
 & and finished sequence \\
\end{tabular}
\end{center}

The course also named Mexican contributions:
Francisco Bolívar Zapata (recombinant insulin
and somatostatin work in Boyer's group), Luis
Herrera-Estrella (transgenic plants), and the
Cry toxins of \emph{Bacillus thuringiensis}
studied by Alejandra Bravo and Mario Soberón.

\textbf{Cloning.}
A DNA fragment cut with a restriction enzyme
is ligated into a vector cut with the same
enzyme; the recombinant plasmid is introduced
into bacteria, transformants are selected
(antibiotic resistance), and recombinants are
screened (blue-white screening, restriction
analysis, sequencing) \cite[Section
3-5]{voet2013fundamentals}. A cDNA copy of an
mRNA, made by reverse transcriptase, lacks
introns and can be expressed in bacteria
\cite[Chapter 10]{alberts2014essential}.
Products include insulin, growth hormone,
erythropoietin and clotting factors
\cite[Section 3-5]{voet2013fundamentals}.

\textbf{PCR.}
The polymerase chain reaction amplifies a
segment between two primers through cycles of
denaturation (94--95 $^\circ$C), annealing of
the primers (typically 50--65 $^\circ$C) and
extension by a heat-stable DNA polymerase (72
$^\circ$C); each cycle roughly doubles the
target \cite[Section
3-5]{voet2013fundamentals}. A reaction
contains template, two primers, the four
dNTPs, $\mathrm{Mg^{2+}}$ and polymerase.
Course data on polymerases: \emph{Taq} has no
3'-5' proofreading exonuclease; Pwo, Pfu and
Pfx do, with fidelities about 12, 30 and 48
times that of \emph{Taq} and longer maximal
products. Primers are usually 18--30
nucleotides with about 50\% G+C, without
self-complementarity (especially at the 3'
ends, to avoid primer dimers). A rough
melting temperature for short primers is the
Wallace rule
$$
T_m = 4(\mathrm{G}+\mathrm{C}) +
2(\mathrm{A}+\mathrm{T}),
$$
in $^\circ$C, counting bases, with annealing
about 5 $^\circ$C below $T_m$.

\textbf{Sequencing.}
In Sanger's chain-terminator method, DNA
polymerase copies a single-stranded template
from a primer in the presence of the four
dNTPs and a small amount of
2',3'-dideoxynucleotides (ddNTPs), which lack
the 3'-OH and stop the chain. Originally four
reactions (one per ddNTP, radioactive primer)
were run side by side on a polyacrylamide
gel; today the four ddNTPs carry different
fluorescent dyes and the products are read in
a capillary, up to about 1000 nucleotides per
read. Newer methods (pyrosequencing,
detection of released protons) sequence
millions of fragments in parallel
\cite[Section 3-4]{voet2013fundamentals}.

\textbf{Gene expression.}
The transcriptome can be surveyed by Northern
blots (one RNA at a time), RT-PCR (reverse
transcription followed by PCR; quantitative
in real time), DNA microarrays (labelled cDNA
hybridized to thousands of probes on a slide)
and RNA-seq (sequencing of all cDNAs,
counting reads per gene) \cite[Chapter
10]{alberts2014essential}.

\section{Detecting genetically modified soy}
\label{section-mcb-gmo}

Most transgenic crops carry an
insect-resistance gene (a Cry toxin of
\emph{B. thuringiensis}: Bt crops) or
herbicide tolerance. Transgenes are built
with plant regulatory elements, and two of
the most common are the constitutive
cauliflower mosaic virus 35S promoter (strong
in dicots) and the terminator of the
\emph{Agrobacterium tumefaciens} nopaline
synthase gene (NOS), which ends transcription
and directs polyadenylation. Detecting either
by PCR therefore screens for most GM events.

The course manual's protocol: grind 100 g of
soy in saline; extract DNA with CTAB buffer
at 65 $^\circ$C, proteinase K, chloroform,
CTAB precipitation, a second chloroform step,
isopropanol precipitation and a 75\% ethanol
wash; check purity by $A_{260}/A_{280}$ (an
$A_{260}$ of 1.0 corresponds to about 50
$\mu$g/mL of double-stranded DNA); amplify
with primers for 35S (196 bp product) and NOS
(180 bp) in 30 cycles of 94, 54 and 72
$^\circ$C; run on agarose.

In the course, soy sprouts bought in an
ordinary Mexico City shop were tested.
Daniel's report covers the extraction and
gel: two sharp high-molecular-weight bands
(estimated at about 10 kb and 6 kb, since the
gel was not run to completion) and a smear
near the bottom, which he attributed, after
weighing protein contamination and
overloading, to degradation by nucleases. The
extraction was judged good enough for the PCR
step. His introduction put the test in
context: by 2013 about 79 of 107 million
hectares of soy worldwide were transgenic
(ISAAA figures, not checked), and in the
United States the share rose from 17\% in
1997 to 94\% in 2015.

\bibliography{my}
\bibliographystyle{amsalpha}

\end{document}
